\chapter{Introdução} \label{cap:intro}

A qualidade da educação básica é um dos principais determinantes do desenvolvimento humano e econômico de longo prazo \cite{hanushek2012}. No Brasil, o Índice de Desenvolvimento da Educação Básica (IDEB) monitora o ensino público combinando fluxo escolar e proficiência em língua portuguesa e matemática numa medida comparável entre municípios, estados e redes de ensino \cite{fernandes2007}. Apesar dos avanços das últimas duas décadas, as disparidades regionais persistem: estados do Norte e do Nordeste ficam sistematicamente abaixo da média nacional, reflexo de desigualdades estruturais de renda, infraestrutura e serviços públicos \cite{inep2023}.

O Maranhão está entre os casos mais críticos. Com um dos menores IDHMs do país e indicadores educacionais historicamente abaixo da média nacional, o estado concentra municípios em alta vulnerabilidade socioeconômica, onde os obstáculos para melhorar o desempenho escolar superam a capacidade de gestão local \cite{ibge2010, inep2023}. Isso coloca uma pergunta direta: o que diferencia, dentro de um mesmo patamar de vulnerabilidade, os municípios que sobem no IDEB daqueles que ficam estagnados?

A literatura internacional examinou amplamente a relação entre condições socioeconômicas, infraestrutura escolar, qualificação docente e o desempenho dos estudantes \cite{rivkin2005, coleman1966}. No Brasil, estudos com dados do IDEB e do Censo Escolar apontaram formação dos professores, tamanho das turmas e infraestrutura como variáveis relevantes \cite{vasconcelos2021, soares2013}. Análises voltadas especificamente para o Maranhão que isolem fatores intracontextuais (aqueles que operam mesmo quando o contexto socioeconômico é semelhante) são escassas.

Métodos de aprendizado de máquina têm sido aplicados a dados educacionais com resultados promissores, especialmente por sua capacidade de lidar com bases administrativas grandes e heterogêneas \cite{baker2010, romero2020}. Algoritmos de agrupamento identificam estruturas nos dados sem categorias predefinidas, o que os torna úteis para segmentar contextos antes de comparações \cite{faceli2021}. Quando combinados com medidas de importância de variáveis e testes não paramétricos, essa abordagem permite identificar os fatores que mais diferenciam grupos de desempenho.

Este trabalho aplica essa lógica aos 217 municípios maranhenses com dados disponíveis no IDEB 2023 (Anos Iniciais do Ensino Fundamental, rede municipal). No primeiro estágio, os municípios são agrupados por contexto socioeconômico via K-Means, usando renda per capita e IDHM. No segundo, dentro de cada grupo, um Random Forest Classifier combinado com o teste de Mann-Whitney identifica quais características escolares e socioeconômicas separam municípios de alto e baixo desempenho. As variáveis analisadas vêm do Censo Escolar 2023: infraestrutura (internet, biblioteca, laboratório de informática, saneamento), qualificação docente e indicadores de fluxo escolar.

O trabalho faz três contribuições principais. Primeiro, propõe e aplica ao contexto específico do Maranhão uma abordagem em dois estágios (segmentação socioeconômica seguida de análise intragrupo). Segundo, identifica os fatores escolares associados ao desempenho superior no IDEB entre municípios vulneráveis. Terceiro, disponibiliza um pipeline reprodutível para análises similares em outras unidades federativas. Os resultados têm uso direto para gestores e formuladores de políticas que precisam decidir onde concentrar recursos.

O restante do artigo está organizado da seguinte forma: o Capítulo~\ref{cap:related} apresenta os trabalhos relacionados; o Capítulo~\ref{cap:methods} descreve os materiais e métodos; o Capítulo~\ref{cap:experiments} apresenta os experimentos e resultados; o Capítulo~\ref{cap:conclusions} discute conclusões, limitações e direções futuras.